\honest{Evaluability (And Cheap Holiday Shopping)}{Eliezer Yudkowsky}{2007}

A reader asks, I pretend, how to be seen as generous without spending much. The answer is yes: according to Hsee, a \$45 scarf is seen as more generous than a \$55 coat. \nb{Hsee's respondents were told that the scarf was near the top of a \$5 to \$50 range and the coat near the bottom of a \$50 to \$500 range. That is what made the scarf look generous.}

``This is a special case of a more general phenomenon,'' I say. \nb{Hsee's gift study did not compare the two gifts side by side. A 2023 replication did, and most people shown both still chose the scarf.} In Hsee's dictionary study, students would pay \$24 for a like-new dictionary with 10,000 entries and \$20 for a torn one with 20,000, when each saw only one. Seen side by side, the torn one got \$27 and the like-new one \$19. \nb{The gap when each was seen alone was not significant by itself; the reversal as a whole was.} A number of entries means little alone, while a torn cover is plainly bad.

Slovic and colleagues found that a 7/36 chance to win \$9 was rated more attractive once a 29/36 chance of losing 5 cents was added, and that far more students then preferred it to a sure \$2. ``You can make a gamble more attractive by adding a strict loss!'' ``Of course,'' I add, ``it only works if the subjects don't see the two gambles side-by-side.'' \nb{The studies did not test that. It is a prediction, not a finding.}

In Hsee's ice cream study, people seeing one serving would pay more for 7 ounces overfilling a small cup than for 8 ounces in a large one; seeing both, they reversed. \nb{The post gives the amounts only in a picture.}

So, for your shopping: ``Decide how much money you want to spend on impressing the recipient, then find the most worthless object which costs that amount.'' The cheaper the class, the more expensive the item looks. I would like to know what you bought.
